% Chapter 10, generated by make_chapters.py from papers/Ch10_Bio_19_SharpP_Hardness/anccft32.tex.
% Source paper: Bio 19
\chapter[Double-stranded reconstruction is \#P-hard]{\texorpdfstring{Why there is no closed formula for $\#\DS$: double-stranded reconstruction is $\sP$-hard}{Why there is no closed formula for #DS: double-stranded reconstruction is #P-hard}}\label{chap:ch10}
\begin{chaptersummary}
Single-stranded reconstruction counts are given by one determinant, the BEST theorem. The
double-stranded count $\#\DS(S,k)$ of Chapter~\ref{chap:ch07} is a sum of such determinants over the points of a lattice
in a box. Chapters~\ref{chap:ch08} and~\ref{chap:ch09} evaluated it for $\phi$X174 by frontier elimination and by visiting
$1.24\cdot10^9$ points. It is natural to ask for a formula that avoids the sum: a Gauss sum over a
discriminant form, a Pfaffian, anything computable in polynomial time. We show that no such formula
exists unless $\mathrm{FP}=\sP$. From every connected loopless graph $G$ of maximum degree~$4$ with even
degrees, we build in polynomial time a circular DNA sequence $S$ of length linear in $|E(G)|$ and an
even $k$ such that
\[
\#\DS(S,k)=2\,\varepsilon(G),
\]
where $\varepsilon(G)$ is the number of Eulerian circuits of $G$. The number of Eulerian circuits is
$\sP$-complete for $4$-regular graphs (Ge and \v Stefankovi\v c, after Brightwell and Winkler). The
mechanism is simple: palindromic $k$-mers make the reverse-complement involution act trivially on
vertices, so that choosing a strand for each $\rho$-orbit is the same as orienting an edge. Two
corollaries follow. The number of box points, Chapter~\ref{chap:ch08}'s first column, is the number of Eulerian
orientations of $G$, which is $\sP$-complete as well. And on this class the generation conjecture of
Chapter~\ref{chap:ch11} holds, by Kotzig's theorem. All identities are verified exactly on $23$ graphs by
\texttt{bio19\_verify.py}. The lattice sum, and the frontier and enumeration methods used to evaluate
it, are therefore not an artefact of the approach; they are what the problem requires.
\end{chaptersummary}

\section{The question}\label{ch10:sec:q}

For one strand the count of sequences with a given $(k+1)$-mer spectrum is a determinant times
factorials \cite{BEST}. For two strands, Theorem~\ref{ch07:thm:lattice} gives
\[
\#\DS(S,k)=\sum_{f\in\Lm,\ |f|\le m,\ f\equiv m\ (2)}\Nseq\bigl(\tfrac{m+f}2\bigr)
\]
Chapter~\ref{chap:ch08} counted the $2\,481\,687\,900$ points of the box for $\phi$X174 at $k=10$, and Chapter~\ref{chap:ch09}
evaluated the sum, $31\,925\,753\,246\,212$. Each step needed an exponential-size
object: a frontier of width $26$, or a visit to every point. A proposal for this chapter asked for a
``meshless'' formula instead. It would be a Gauss sum built on the symmetric linking form of the double
cover's sandpile group (Remark~\ref{ch06:rem:pairing}), or a Pfaffian in a Clifford algebra, evaluated in time polynomial in the
number of vertices. We prove that such a formula cannot exist unless every problem in $\sP$ is solvable
in polynomial time.

\section{Eulerian circuits and their orientations}\label{ch10:sec:euler}

Let $G$ be a connected loopless multigraph whose degrees are all even, with at least three edges. An
\emph{Eulerian circuit} is a cyclic sequence of its edges, traversing each once, considered up to
rotation and reversal. Let $\varepsilon(G)$ be their number. An \emph{Eulerian orientation} directs
each edge so that every vertex has equal in- and out-degree. For a balanced digraph $O$ with labelled
arcs, $\ec(O)=\tau(O)\prod_v(d^+_v-1)!$ counts its Eulerian circuits up to rotation \cite{BEST}.

\begin{lemma}\label{ch10:lem:orient}
$\displaystyle\sum_{O\ \text{Eulerian orientation of }G}\ec(O)=2\,\varepsilon(G)$.
\end{lemma}

\begin{proof}
An Eulerian circuit of $G$, traversed in one of its two directions, orients every edge. It is balanced at
every vertex, since each visit enters and leaves once, and it gives a directed Eulerian circuit of that
orientation. Conversely, a directed Eulerian circuit of an orientation, forgetting the arrows, is a
traversed Eulerian circuit of $G$. So the left side counts pairs (circuit, direction). The two directions
of a circuit are different directed circuits, because they traverse the first edge oppositely.
\end{proof}

\begin{theorem}[\cite{BW,GS12}]\label{ch10:thm:sharpP}
Computing $\varepsilon(G)$ is $\sP$-complete, even for $4$-regular graphs.
\end{theorem}

Counting the Eulerian orientations of a graph is also $\sP$-complete \cite{MW96}.

\section{Genomes whose double cover is a given graph}\label{ch10:sec:real}

Let $G$ have vertex set $\{1,\dots,n\}$, $m\ge3$ edges, and all degrees in $\{2,4\}$. Fix an even word length
$k=2h$ and a spacer length $\ell$. The construction has three steps.
\begin{enumerate}
\item To each vertex $x$ assign a palindrome $P_x=u_x\,\rc(u_x)$ with $|u_x|=h$.
\item Compute an Eulerian circuit $x_0\xrightarrow{e_1}x_1\xrightarrow{e_2}\cdots\xrightarrow{e_m}x_m=x_0$ of $G$.
Spell
\[
S=P_{x_0}\,\Sigma_1\,P_{x_1}\,\Sigma_2\cdots P_{x_{m-1}}\,\Sigma_m\qquad(\text{circular}),
\]
with one spacer $\Sigma_i$ of length $\ell$ per traversal.
\item Assign the boundary letters of the spacers. $P_x$ occurs $\deg(x)$ times in $S$ and $\rc S$
together. The letters that follow it are the first letters of the outgoing spacers and the complements
of the last letters of the incoming ones. At the $j$-th visit to $x$ ($j=0,1$), let the outgoing spacer
begin with $\mathtt{AC}[j]$ and the incoming one end with $\mathtt{CA}[j]$. The letters following the
occurrences of $P_x$ are then $\mathtt A,\mathtt G$ for a vertex of degree~$2$ and
$\mathtt A,\mathtt C,\mathtt G,\mathtt T$ for a vertex of degree~$4$, pairwise different. The interior
letters of the spacers are free.
\end{enumerate}

\begin{lemma}[realization]\label{ch10:lem:real}
Suppose that in $S\cup\rc S$ the only $k$-mers occurring more than once are the $P_x$, the $P_x$ are
distinct, and no $P_x\Sigma P_y$ is its own reverse complement. Then the double cover $D_k(S)$ has the
following properties:
\begin{enumerate}
\item its vertices are the $n$ palindromes $P_x$, so $\bar v=v$ for every vertex;
\item its contents are the words $P_{x_{i-1}}\Sigma_iP_{x_i}$ and their reverse complements, one
$\rho$-orbit of multiplicity $m=1$ per edge of $G$, with no self-complementary content;
\item the quotient of $D_k(S)$ by $\rho$ is $G$.
\end{enumerate}
\end{lemma}

\begin{proof}
The vertices of $D_k$ are the repeated $k$-mers (Definition~\ref{ch07:def:cover}); $P_x$ occurs $\deg(x)\ge2$ times. The contents
are the stretches between consecutive vertex occurrences, i.e.\ the $P\Sigma P$ words of $S$ and their
reverse complements in $\rc S$. Each occurs once, so it has multiplicity $1$. The orbit of the $i$-th one
joins $P_{x_{i-1}}$ and $P_{x_i}$, which are fixed by $\rho$.
\end{proof}

\begin{lemma}[the hypotheses can be met in polynomial time]\label{ch10:lem:derand}
There are $h,\ell=O(\log m)$ and a deterministic algorithm, polynomial in $m$, that chooses the words $u_x$
and the spacer interiors so that the hypotheses of Lemma~\ref{ch10:lem:real} hold.
\end{lemma}

\begin{proof}[Proof sketch]
Choose the free letters uniformly at random, with $\ell\ge k$. A \emph{bad event} is one of three things:
two occurrences in $S\cup\rc S$ of equal $k$-mers that are not both occurrences of one $P_x$; $P_x=P_y$
for some $x\ne y$; or a self-complementary $P\Sigma P$. Since $\ell\ge k$, a window of length $k$ is an
occurrence of some $P_x$, or it lies inside one spacer, or it straddles one junction between a
palindrome and a spacer. In the last two cases it contains at least $h-2$ free letters from a single
source: half of one palindrome, or the interior of one spacer.

Two windows that differ only in fixed letters cannot be equal, by the boundary assignment. This is the
reason for it: without it the $(k+1)$-mers $P_x\mathtt c$ would repeat with probability
$1-4!/4^4\approx0.91$ at a vertex of degree~$4$. Otherwise, equality of two windows is a system of
equations $a_i=b_i$ or $a_i=\overline{b_i}$ between letters. The complement case arises when one
window lies in $\rc S$. An equation $a=\bar a$ has no solution over $\{\mathtt A,\mathtt C,\mathtt G,\mathtt T\}$,
and the remaining equations link the free letters in paths and cycles. So the system fixes at least half
of them, and its probability is at most $4^{-(h-2)/2}$. There are at most $(2m(k+\ell))^2$ pairs of
windows. With $h\ge 2\log_4\bigl(8m^2(k+\ell)^2\bigr)+2$ the expected number of bad events is below $1/2$.

The conditional probability of each bad event, given some letters, is computed exactly from the same
system of equations. So the method of conditional expectations fixes the letters one at a time,
keeping the expected number of bad events below $1$. This is deterministic and polynomial.
\end{proof}

The verification script uses random interiors and checks the hypotheses, retrying on failure. With
$k=12$ and $\ell=18$ it never needed a retry on the graphs below.

\section{The count}\label{ch10:sec:count}

\begin{theorem}\label{ch10:thm:main}
For $S$ and $k$ as in Section~\ref{ch10:sec:real},
\[
\#\DS(S,k)=2\,\varepsilon(G),
\]
and the box points of Theorem~\ref{ch07:thm:lattice} are in bijection with the Eulerian orientations of $G$.
\end{theorem}

\begin{proof}
Every orbit has $m=1$. So a point of the box chooses, for each orbit $\{a,\rho a\}$, exactly one of $a$
and $\rho a$. Since all vertices are palindromic, $a\colon P_x\to P_y$ and $\rho a\colon P_y\to P_x$
are the two orientations of the edge $\{x,y\}$ of $G$. A point is a circulation iff the chosen
orientation is balanced, i.e.\ Eulerian. At it, $c(a)\in\{0,1\}$, so the factorials of
$\Nseq=\tau\prod(d^+-1)!/\prod c(a)!$ disappear and $\Nseq=\ec(O)$ (Theorem~\ref{ch07:thm:lattice}). Summing gives
Lemma~\ref{ch10:lem:orient}.
\end{proof}

\begin{corollary}\label{ch10:cor:hard}
Given a circular DNA sequence $S$ and an even $k$, computing $\#\DS(S,k)$ is $\sP$-hard, and so is
computing the number of points of Chapter~\ref{chap:ch07}'s box. In particular, no algorithm, closed formula or
``meshless'' expression evaluates $\#\DS$ in time polynomial in $|S|$ unless $\mathrm{FP}=\sP$.
\end{corollary}

\begin{proof}
The map $G\mapsto(S,k)$ is polynomial and $|S|=m(k+\ell)=O(m\log m)$. Divide by $2$ and apply
Theorem~\ref{ch10:thm:sharpP}. For the box, use \cite{MW96}.
\end{proof}

\begin{table}[t]
\centering\small
\begin{tabular}{lrrrrrr}
\toprule
graph & $n$ & $|E|$ & $|S|$ & $\varepsilon(G)$ & Eulerian orientations & $\#\DS(S,12)$\\
\midrule
$4$-regular & 3 & 6 & 180 & 16 & 10 & 32\\
$4$-regular & 4 & 8 & 240 & 44 & 16 & 88\\
$4$-regular & 5 & 10 & 300 & 112 & 28 & 224\\
$4$-regular & 6 & 12 & 360 & 328 & 42 & 656\\
$4$-regular, simple & 5 & 10 & 300 & 132 & 24 & 264\\
$4$-regular, simple & 6 & 12 & 360 & 372 & 38 & 744\\
$4$-regular, simple & 7 & 14 & 420 & 1052 & 60 & 2104\\
degrees $2,4$ & 6 & 10 & 300 & 36 & 18 & 72\\
degrees $2,4$ & 4 & 6 & 180 & 6 & 6 & 12\\
\bottomrule
\end{tabular}
\caption{A selection of the $23$ graphs of checks (E), (R), (D), (K). $\varepsilon(G)$ is found by exhaustive
enumeration of trails. $\#\DS$ comes from Chapter~\ref{chap:ch07}'s lattice sum on the double cover built from the
definition, and it equals $2\varepsilon(G)$ in every case. So does the closure of $S$ under Chapter~\ref{chap:ch11}'s
moves. The box points equal the Eulerian orientations in every case.}\label{ch10:tab:res}
\end{table}

Table~\ref{ch10:tab:res} shows a selection of the verified cases. In every case, the double cover of the
constructed genome, built from the definition, has exactly $n$ palindromic vertices, $|E|$ orbits of
multiplicity $1$ and no self-complementary content. In check~(S), $4$-regular graphs with $n=4,\dots,8$
give genomes of $240$--$480$ bp, linear in $|E|$. On them $\#\DS$ grows from $80$ to $6\,008$.

\section{Two by-products}\label{ch10:sec:by}

\paragraph{Chapter~\ref{chap:ch11}'s conjecture on this class.} Conjecture~\ref{ch11:conj:gen} of Chapter~\ref{chap:ch11} states that transpositions, inversions
at inverted repeats, and reverse complementation connect $\DS(S,k)$. Here every vertex is its
own complement. So an inversion between two visits of $P_x$ (a flip in the sense of
Definition~\ref{ch07:def:flip}) reverses the closed stretch of the circuit between them, which is a
\emph{$\kappa$-transformation}. Kotzig proved that the $\kappa$-transformations
connect all Eulerian circuits of a connected graph \cite{Kot68,Fle90}. Reversal of the whole circuit
supplies the direction.

\begin{corollary}
Conjecture~\ref{ch11:conj:gen} holds for every genome of Section~\ref{ch10:sec:real}.
\end{corollary}

Check~(K) confirms it: the closure of $S$ under inversions and reverse complementation alone is all of
$\DS(S,k)$ in all $23$ cases.

\paragraph{Both counting columns are hard.} Chapter~\ref{chap:ch08} counted the points of the box ($2\,481\,687\,900$ for
$\phi$X174 at $k=10$) and Chapter~\ref{chap:ch09} the reconstructions. By Theorem~\ref{ch10:thm:main} the first count is the
number of Eulerian orientations of a graph, and the second twice its number of Eulerian circuits. Both are
$\sP$-complete in general.

\section{What this does and does not exclude}\label{ch10:sec:scope}

\begin{itemize}
\item \emph{Excluded:} any method polynomial in $|S|$ that is exact on all inputs. That includes a
Pfaffian, a determinant, or a Gauss sum with polynomially many polynomially computable terms. Unless
$\mathrm{FP}=\sP$, none of them exists. The proposed Gauss-sum expansion
$\#\DS=|K|^{-1}\sum_\gamma\mathcal G(\gamma)\Phi(\gamma)$ fails for a second, independent reason. The
summation domain is the box, not a group, so Poisson summation over the finite group $K(D_k)$ does not
collapse it. Written as a Fourier expansion on $K$, the $\Phi(\gamma)$ carry the entire difficulty.
\item \emph{Not excluded:} algorithms exponential in a structural parameter. Chapter~\ref{chap:ch08}'s frontier
elimination is exponential in the frontier width and linear otherwise. Genomes of bounded width are
tractable. $\phi$X174 at $k=10$ has width $26$, which is why it was feasible.
\item \emph{Not excluded:} efficient approximation. Whether $\#\DS$ admits a fully polynomial randomized
approximation scheme is open; so is the corresponding question for undirected Eulerian circuits.
\item \emph{Odd $k$.} The construction needs palindromic $k$-mers, which exist only for even $k$. For odd
$k$ no vertex is its own complement, and the quotient of $D_k$ is the signed graph of Chapter~\ref{chap:ch06}.
The strand choices are then orientations of a signed graph, and the count is one of Eulerian circuits of a
bidirected graph. We expect hardness there too, but do not prove it.
\end{itemize}

\section*{Ancillary file}
\texttt{bio19\_verify.py} (standard library; imports \texttt{ds\_fermion.py} of Chapter~\ref{chap:ch09} and
\texttt{bio17\_verify.py} of Chapter~\ref{chap:ch11}) runs checks (E), (R), (D), (K) and (S) in two seconds. It writes
its transcript to \texttt{bio19\_output.txt}.
